\csection{Введение}

Современные \eng{e-commerce} проекты используют несколько каналов привлечения пользователей: органический поиск, контекстную рекламу, социальные сети, прямые переходы, \eng{referral}-трафик и \eng{email}-рассылки. Каждый канал формирует собственный поток данных о посещениях, событиях, заказах, возвратах и выручке. При отсутствии единой аналитической среды эти данные часто хранятся в разрозненных таблицах, отчетах рекламных кабинетов и внутренних учетных системах, что усложняет оценку реальной эффективности продвижения.

Особенно важной задачей для \eng{SEO}-специалиста является определение источников трафика, которые приводят не только к переходам, но и к целевым действиям: просмотру товара, добавлению в корзину, оформлению заказа и покупке. Ручная обработка таких данных требует значительных временных затрат, повышает вероятность ошибок и затрудняет подготовку регулярной отчетности для клиента.

Актуальность темы обусловлена тем, что оценка эффективности продвижения интернет-магазина не может ограничиваться только количеством посещений сайта. Для принятия управленческих решений необходимо понимать, как пользователь проходит путь от первого перехода до покупки, какие этапы воронки вызывают потери, какие источники формируют заказы с высокой выручкой, а какие приводят только к нецелевому трафику. При этом данные о действиях пользователей, заказах, возвратах и расходах обычно имеют разную структуру, разную периодичность обновления и разные правила интерпретации.

В условиях учебного дипломного проекта важно показать не только интерфейс аналитической панели, но и саму логику обработки данных. Разрабатываемое программное средство должно демонстрировать полный цикл работы: от регистрации заявки клиента до просмотра аналитического отчета. Поэтому в проекте рассматриваются роли клиента, \eng{SEO}-специалиста и администратора, а также процессы ввода данных, проверки их корректности, расчета показателей и формирования рекомендаций. Такой подход позволяет представить систему как законченное программное средство, а не как набор отдельных экранов.

Отдельное значение имеет то, что проект является образовательным и не предполагает фактической коммерческой эксплуатации. Внешний интернет-магазин и сторонние сервисы не рассматриваются как отдельные акторы системы. Вместо этого используется ввод данных в табличном виде и генерация демонстрационного массива данных администратором. Данное решение позволяет сосредоточиться на проектировании сервиса сквозной аналитики, модели данных, интерфейсе и серверной логике без необходимости разрабатывать внешнюю торговую площадку.

Целью дипломного проекта является разработка программного средства реализации сервиса сквозной аналитики для \eng{e-commerce}, предназначенного для автоматизации подключения и сопровождения аналитики интернет-магазинов. Разрабатываемая система должна обеспечивать регистрацию заявок клиентов, ведение проектов, ввод бизнес-данных в табличном виде, генерацию демонстрационных данных, расчет аналитических показателей, визуализацию \eng{dashboard}, формирование рекомендаций и подготовку отчетов.

Для достижения поставленной цели необходимо решить следующие задачи: изучить предметную область сквозной аналитики интернет-магазинов; проанализировать существующие программные решения; выделить недостатки ручной обработки данных; формализовать бизнес-процессы текущего и предлагаемого состояния; определить функциональные и нефункциональные требования; спроектировать архитектуру клиентской и серверной частей; разработать модель данных; описать механизмы разграничения прав доступа; подготовить пользовательский интерфейс для основных ролей системы.

Объектом исследования является процесс подключения и сопровождения сквозной аналитики интернет-магазина. Предметом исследования являются методы и программные средства сбора, проверки, агрегации и визуального представления аналитических данных по источникам трафика, событиям, заказам, выручке и возвратам.

Практическая значимость работы заключается в том, что результат дипломного проекта может использоваться как демонстрационная система для анализа принципов сквозной аналитики. Программное средство позволяет показать, каким образом данные из разных таблиц объединяются в единую аналитическую картину, как рассчитываются показатели эффективности и как результаты отображаются пользователю в виде \eng{dashboard}, отчетов и рекомендаций.

В рамках дипломного проекта рассматриваются предметная область сквозной аналитики, существующие программные решения, бизнес-процессы текущего и предлагаемого состояния, требования к программному средству, архитектура \eng{web}-системы, модель данных, пользовательский интерфейс, механизмы информационной безопасности и условная оценка экономической эффективности внедрения.
